\naslov{Centralni limitni izrek}

Naj bo $S_n \sim \operatorname{Bin}(n, \pol)$, da velja $P(S_n = k) =
\binom{n}{k} 2^{-n}$.
Poglejmo si razmerje
\[
  \frac{P(S_n = k+1)}{P(S_n = k)}
  = \frac{\binom{n}{k+1}}{\binom{n}{k}}
  = \frac{n-k}{k+1}.
\]
Za sedaj se omejimo na sode $n = 2m$.
Potem je $P(S_n = k)$ največja za $k=m$, in za $d > 0$ dobimo
\begin{align*}
  \frac{P(S_{2m} = m+d)}{P(S_{2m} = m)}
  &= \frac{P(S_{2m} = m+d)}{P(S_{2m} = m+d-1)} \frac{P(S_{2m} = m+d-1)}{P(S_{2m}
	= m+d-2)} \cdots \frac{P(S_{2m} = m+1)}{P(S_{2m} = m)} \\
  &= \frac{2m - m - d +1}{m+d} \cdots \frac{2m-m}{m+1} \\
  &= \frac{m - d + 1}{m - d} \cdots \frac{m}{m+1} \\
  &= \frac{1 + \frac{1-d}{m}}{1 + \frac{d}{m}} \cdots \frac{1}{1 + \frac{1}{m}} \\
  &= \frac{1 - \frac{d-1}{m}}{1 + \frac{1}{m}} \cdots \frac{1}{1 + \frac{d}{m}}.
\end{align*}
Če je $d$ dovolj majhen, je to približno enako
\begin{align*}
  \frac{P(S_{2m} = m+d)}{P(S_{2m} = m)}
  &\approx \frac{e^{- \frac{d-1}{m}} e^{- \frac{d-2}{m}}\cdots e^0}{e^{\frac{1}{m}} e^{\frac{2}{m}} \cdots e^{\frac{d}{m}}} \\
  &= \exp\left( - \frac{1}{m} - \frac{2}{m} - \cdots - \frac{d}{m} - \frac{d-1}{m} - \cdots - \frac{1}{m} \right) \\
  &= \exp\left( -\frac{d^2}{m} \right).
\end{align*}
Pri $d < 0$ je porazdelitev simetrična, torej dobimo enak rezultat.
Izkaže se
\[
  \lim_{m \to \infty} \sum_{d = -\infty}^\infty \abs{ \frac{P(S_{2m} =
	  m+d)}{P(S_{2m} = m)} - e^{-d^2 / m} } = 0.
\]
Vrsta
\[
  \sum_{d = -\infty}^\infty \frac{P(S_{2m} = m+d)}{P(S_{2m} = m)} =
  \inv{P(S_{2m} = m)}
\]
je približno enaka
\[
  \sum_{d = -\infty}^\infty e^{-d^2/m} \approx \int_{-\infty}^\infty e^{-x^2 /
	m} dx = \sqrt{m\pi},
\]
torej je verjetnost $P(S_{2m} = m) \approx \inv{\sqrt{m \pi}}$.
Upoštevaje zgornjo formulo potem dobimo
\[
  P(S_{2m} = m+d) \approx \inv{\sqrt{m \pi}} e^{- \frac{d^2}{m}}
\]
oziroma
\[
  P(S_n = k) \approx \sqrt{\frac{2}{n \pi}} \exp\left( - \frac{2
	  (k-\nicefrac{n}{2})^2}{n} \right).
\]
Vse to je res tudi za lihe $n$.
Rezultatu pravimo \pojem{de Moivreova lokalna formula}.

\vprasanje{Izpelji de Moivreovo lokalno formulo.}

Za $a, b \in \Z$ lahko izpeljemo
\begin{align*}
  P(a \le S_n \le b)
  &= \sum_{k=a}^b P(S_n = k) \\
  &\approx \sqrt{\frac{2}{n \pi}} \sum_{k=a}^b \exp\left( - \frac{2 (k - \frac{m}{2})^2}{n} \right) \\
  &\approx \sqrt{\frac{2}{n \pi}} \sum_{k=a}^b \int_{k - \pol}^{k + \pol} \exp\left( - \frac{2 (x - \frac{m}{2})^2}{n} \right) dx \\
  &= \sqrt{\frac{2}{n \pi}} \int_{a}^{b} \exp\left( - \frac{2 (x - \frac{m}{2})^2}{n} \right) dx,
\end{align*}
kar se da razširiti za poljubna $a < b$.
Dobljeni približek za verjetnost $P(S_n = k)$ je gostota normalne porazdelitve
$N(\frac{n}{2}, \frac{n}{4})$ v točki $k$.
Velja $E(S_n) = \frac{n}{2}$ in $\var(S_n) = \frac{n}{4}$.
Binomska porazdelitev, s katero smo začeli, je enaka vsoti $n$ neodvisnih
Bernoullijevih slučajnih spremenljivk, ki je torej porazdeljena približno
normalno.

\begin{izrek}[centralni limitni izrek]
  Naj bodo $X_1, X_2, \ldots$ neodvisne in enako porazdeljene slučajne
  spremenljivke s končnim drugim momentom.
  Označimo $\mu_1 = E(X_i)$ in $\sigma_1^2 = \var(X_i)$.
  Za $S_n = X_1 + \cdots + X_n$ veljata naslednji točki:
  \begin{itemize}
  \item Če definiramo $W_n = \frac{S_n - n \mu_1}{\sigma_1 \sqrt{n}}$, za vsak
	$w \in \R$ velja
	\[
	  \lim_{n \to \infty} P(W_n \le w) = \phi(w).
	\]
  \item Limita
	\[
	  \lim_{n \to \infty} \sup_{w \in \R} \abs{P(W_n < w) - \phi(w)}
	  = \lim_{n \to \infty} \sup_{x \in R} \abs{P(S_n < x) - \phi\!\left(
		  \frac{x - n \mu_1}{\sigma_1 \sqrt{n}} \right)}
	  = 0.
	\]
  \end{itemize}
\end{izrek}

\vprasanje{Formuliraj centralni limitni izrek.}

Dokažemo lahko tudi drugačno formulacijo:
za vsako zvezno in omejeno funkcijo $h: \R \to \R$ z največ kvadratno rastjo
(tj. obstaja $M$, da je $\abs{h(x)} \le M(1+x)^2$) velja
\[
  \lim_{n \to \infty} E(h(W_n)) = E(h(Z)),
\]
kjer je $Z \sim N(0,1)$.
Funkcijam $h$ pravimo \pojem{poizkusne} ali \pojem{testne} funkcije.

\begin{trditev}[Jensenova neenakost]
  Naj bo $X$ slučajna spremenljivka z vrednostmi na intervalu $I \subseteq \R$
  in $\varphi:I \to \R$ konveksna funkcija.
  Tedaj je $\varphi(E(X)) \le E(\varphi(X))$.
\end{trditev}

Ideja dokaza je, da poiščemo linearno funkcijo, ki se v iskani točki dotika
grafa $\varphi$, in vrednost aproksimiramo s pomočjo nje.
Trditev je pomembna zaradi posledice.

\begin{posledica}
  Če je $p \ge q$, je $E(\abs{X}^q) \le (E(\abs{X}^p))^{q/p}$.
\end{posledica}

\begin{posledica}
  Če obstaja $p$-ti moment in je $q \le p$, obstaja tudi $q$-ti moment.
\end{posledica}

Za nadaljevanje naj bodo $Y_1, \ldots, Y_n$ neodvisne z vsoto $S$ in $Z_1,
\ldots, Z_n$ neodvisne z vsoto $T$.
Privzemimo še, da tretji momenti $Y_k$ in $Z_k$ obstajajo, da je $E(Y_k) =
E(Z_k) = 0$ in da je $\var(Y_k) = \var(Z_k) = \sigma_k^2$.
Tedaj je $E(T) = E(S) = 0$ in
\[
  \var(S) = \var(T) = \sum_{k} \sigma_k^2.
\]
Za primerno $h : \R \to \R$ bomo ocenili razliko $E(h(S)) - E(h(T))$.
Privzeli bomo, da je $h \in \zvodvi{3}$ in $M_3 = \sup_{x \in \R} \abs{h'''(X)}
< \infty$.
Definirajmo kombinirane vsote $V_k = Y_1 + \ldots + Y_k + Z_{k+1} + \ldots +
Z_n$.

Dodatno privzemimo, da so $Y_i$ in $Z_i$ vsi medsebojno neodvisni.
Velja
\[
  E(h(S)) - E(h(T)) = \sum_{k=1}^n E(h(V_k) - h(V_{k-1}))
  = \sum_{k=1}^n E(h(U_k + Y_k) - h(U_k + Z_k))
\]
za $U_k = Y_1 + \ldots + Y_{k-1} + Z_{k+1} + \ldots + Z_n$.
Velja
\[
  h(V_k) = h(U_k) + h'(U_k) Y_k + \pol h''(U_k) Y_k^2 + R_k,
\]
kjer je $R_k$ omejen z $\inv{6} M_3 \abs{Y_k}^3$.
Poleg tega je
\[
  h(V_{k-1}) = h(U_k) + h'(U_k) Z_k + \pol h''(U_k) Z_k^2 + \tilde{R}_k,
\]
kjer je $\tilde{R}_k$ podobno omejena.
Zaradi neodvisnosti $U_k$, $Y_k$ in $Z_k$ potem dobimo
\begin{gather*}
  E(h(V_k)) = E(h(U_k)) + E(h'(U_k)) E(Y_k) + \pol E(h''(U_k)) E(Y_k^2) + E(R_k),
  \\
  E(h(V_{k-1})) = E(h(U_k)) + E(h'(U_k)) E(Z_k) + \pol E(h''(U_k)) E(Z_k^2) + E(\tilde{R}_k).
\end{gather*}
Torej
\[
  \abs{E(h(V_k) - h(V_{k-1}))} \le E(R_k) + E(\tilde{R}_k) \le \inv{6} M_3
  (E(\abs{Y_k}^3) + E(\abs{Z_k}^3)),
\]
iz česar naposled dobimo
\[
  \abs{E(h(S)) - E(h(T))} \le \inv{6} M_3 \sum_{k=1}^n \left( E(\abs{Y_k}^3) +
	E(\abs{Z_k}^3) \right).
\]

\vprasanje{Izpelji oceno za razliko $\abs{E(h(S)) - E(h(T))}$, kjer je $S = Y_1
  + \ldots + Y_n$ in $T = Z_1 + \ldots + Z_n$, in so $Y_i$ enako porazdeljeni
  neodvisni, ter $Z_i$ enako porazdeljeni in neodvisni z $E(Y_i) = E(Z_i) = 0$,
  $\var(Y_i) = \var(Z_i) = \sigma_i^2$.}

Če vzamemo $Z_k \sim N(0, \sigma_k^2)$, je $T \sim N(0, \sigma^2)$.
Tedaj lahko z integralom izračunamo
\[
  E(\abs{Z_k}^3) = \sigma_k^3 \frac{4}{\sqrt{2 \pi}}
  = \frac{4}{\sqrt{2 \pi}} \left( E(Y_k^2) \right)^{3/2}
  \le \frac{4}{\sqrt{2 \pi}} E(\abs{Y_k}^3),
\]
kjer smo v zadnjem koraku uporabili posledico Jensenove neenakosti.
V tem primeru torej lahko ocenimo
\[
  \abs{E(h(S)) - E(h(T))}
  \le \left(\inv{6} + \frac{2}{3 \sqrt{2 \pi}}\right) M_3 \sum_{k=1}^n E(\abs{Y_k}^3).
\]
Tu žal ne moramo vzeti $h(w) = \mathbbm{1}(w \le a)$.
Velja pa naslednji izrek.

\begin{izrek}
  Obstaja taka univerzalna konstanta $C$, da za poljubne neodvisne slučajne
  spremenljivke $Y_1, \ldots, Y_n$, za katere je $E(Y_k) = 0$, velja
  \[
	\sup_{x \in \R} \abs{P(S \le x) - \phi\!\left(\frac{x}{\sigma}\right)}
	= \sup_{w \in \R} \abs{P\!\left(\frac{S}{\sigma} < w\right) - \phi(w)}
	\le \frac{C}{\sigma^3} \sum_{k=1}^n E(\abs{Y_k}^3),
  \]
  pri čemer je $S = Y_1 + \ldots + Y_n$ in $\sigma^2 = \var(S) = \sum_k
  \var(Y_k)$.
\end{izrek}

Naj bodo $X_1, X_2, \ldots$ neodvisne in enako porazdeljene z $E(X_1) = \mu_1$
ter $\var(X_1) = \sigma_1^2$ in $E(\abs{X_1 - \mu_1}^3) = \gamma^3$.
Definiramo $S_n = X_1 + \ldots + X_n$ in
\[
  W_n = \frac{S_n - n \mu_1}{\sigma_1 \sqrt{n}}.
\]
Za fiksen $n$ označimo
\[
  Y_k = \frac{X_k - \mu_1}{\sigma_1 \sqrt{n}},
\]
da je $E(Y_k) = 0$ in $W_n = \sum_k Y_k$.
Velja $\sigma^2 = \var(W_n) = 1$ in
\[
  \sum_{k=1}^n E(\abs{Y_k}^3) = n E(\abs{Y_1}^3) = \frac{\gamma_1^3}{\sigma_1^3
	\sqrt{n}}
  \xrightarrow[{n \to \infty}]{} 0.
\]

\begin{posledica}
  S temi oznakami
  \[
	\sup_{x \in \R} \abs{P(S_n < x) - \phi\!\left( \frac{x - n \mu_1}{\sigma_1
		  \sqrt{n}} \right)}
	= \sup_{w \in \R} \abs{P(W_n < w) - \phi(w)}
	\le \frac{C}{\sqrt{n}} \frac{\gamma_1^3}{\sigma_1^3},
  \]
  za vsako funkcijo $h \in \zvodv{3}{\R}$ pa velja še
  \[
	\abs{E(h(W_n)) - E(h(Z))}
	\le \left( \inv{6} + \frac{2}{3 \sqrt{2 \pi}} \right)
	\frac{\gamma_1^3}{\sigma_1^3} \frac{M_3}{\sqrt{n}}.
  \]
\end{posledica}

\vprasanje{Kaj lahko poveš o oceni razlike porazdelitve vsote od normalne
  porazdelitve?}
